% !TEX root = ../../report.tex

\section{Анализ моделей баз данных}

Для выбора модели хранения данных рассмотрены три группы моделей: дореляционные, реляционные и постреляционные~\cite{db_models}. Сравнение выполняется по признакам, связанным с проектируемой системой: фиксированность схемы, наличие стандартного языка запросов, возможность хранения неструктурированных данных и поддержка целостности.

Дореляционные модели используют иерархические или сетевые связи между записями с жёстко заданной схемой и структурированными записями. Такая организация усложняет изменение связей и выполнение декларативных запросов к данным~\cite{db_models}. Постреляционные модели расширяют реляционный подход средствами хранения составных и полуструктурированных данных. Для текущей задачи такие возможности не являются основными, так как каталог, заказы, пользователи и отзывы имеют заранее определённые атрибуты и связи.

Реляционная модель представляет данные в виде отношений, поддерживает ключи, внешние связи и стандартный язык запросов~\cite{db_models, relation_model}. Эти свойства соответствуют требованиям проектируемой системы: требуется хранить связанные сущности, контролировать целостность данных и разграничивать доступ к ним.

Сравнение рассмотренных групп моделей представлено в таблице~\ref{tbl:model_comparison}.

\begin{table}[H]
    \caption{Сравнение моделей БД}
    \label{tbl:model_comparison}
    \centering
    \begin{tabular}{|l|l|l|l|}
        \hline
        \multicolumn{1}{|c|}{\multirow{2}{*}{\textbf{Критерий}}} & \multicolumn{3}{c|}{\textbf{Модель}} \\ \cline{2-4}
        \multicolumn{1}{|c|}{} & \makecell[l]{Дореля-\\ционная} & \makecell[l]{Реля-\\ционная} & \makecell[l]{Постреля-\\ционная} \\ \hline
        \makecell[l]{Фиксированная\\схема данных} & Да & Да & Частично \\ \hline
        \makecell[l]{Стандартный\\язык запросов} & Нет & Да & Да \\ \hline
        \makecell[l]{Хранение неструкту-\\рированных данных} & Нет & Нет & Да \\ \hline
        \makecell[l]{Обеспечение\\целостности данных} & Частично & Да & Частично \\ \hline
    \end{tabular}
\end{table}

Для проектируемой системы выбрана \textbf{реляционная модель} базы данных. Она соответствует структуре данных онлайн-маркетплейса и позволяет описать связи между товарами, заказами, пользователями, отзывами и категориями средствами ограничений целостности.
